% ============================================================
\section{Conclusiones}
% ============================================================

Del estado del arte se identificaron cinco técnicas de subrogación aplicables a control predictivo: ARX, ARMAX, FOPDT, función de transferencia (TF) y espacio de estados (SS). Sin embargo, el modelo Hammerstein-Wiener, también revisado en la literatura, no es compatible con esquemas de MPC lineal debido a su componente no lineal estático, por lo que solo cuatro familias de modelos resultan integrables de forma directa en el \textit{Model Predictive Control Toolbox}. Dentro de esas cuatro familias, ARX y ARMAX
se implementaron como estructuras separadas, dando lugar a los cinco modelos evaluados en este trabajo.

\vspace{0.25cm}

La conclusión metodológica central del trabajo es que las métricas de ajuste en lazo abierto RECM y $R^2$ no son criterios válidos para la selección de modelos subrogados en aplicaciones de control predictivo. Estas métricas caracterizan la capacidad del modelo para reproducir trayectorias ya conocidas bajo condiciones de identificación, pero no capturan los requisitos que el MPC impone sobre el modelo: representar con fidelidad la dinámica dentro del horizonte de predicción, bajo retroalimentación continua y en  presencia de perturbaciones no contempladas en los datos de entrenamiento. La falta de correlación entre el ranking de ajuste estadístico y el ranking de desempeño en lazo cerrado no es un resultado aislado de este proceso, sino una consecuencia estructural de esa diferencia de objetivos; por lo tanto, la evaluación mediante co-simulación bajo perturbaciones representativas debe considerarse un paso obligatorio en cualquier metodología de subrogación orientada a control. Una segunda conclusión es que la dinámica de cada subsistema determina la arquitectura de modelo más adecuada, y que esta elección no puede generalizarse dentro del mismo proceso. La columna extractiva requiere representar perturbaciones de composición no medidas en la alimentación, lo que exige la capacidad de modelado de ruido correlacionado que aportan los modelos de la familia ARMAX. La columna de recuperación de solvente, en cambio, exhibe una dinámica dominada por efectos integrativos de masa y energía que solo quedan bien representados cuando el modelo incluye variables de estado internas. Imponer una única arquitectura a ambas columnas habría degradado necesariamente el desempeño de control en al menos una de ellas. La propuesta de control híbridodescentralizado un modelo distinto por subsistema no es una solución ad hoc, sino la consecuencia lógica de reconocer que la heterogeneidad dinámica del proceso no puede ser absorbida por una estructura uniforme.

\vspace{0.25cm}

El trabajo también permite concluir que los modelos lineales identificados alrededor del punto nominal de operación son suficientes para garantizar el rechazo de perturbaciones en el rango operativo evaluado. Esto implica que la complejidad adicional de modelos no
lineales no se justifica mientras el proceso opere cerca del punto de diseño, y que la estrategia de linealización local es una aproximación válida para esta clase de procesos de separación bajo las condiciones de perturbación consideradas.Respecto al modelo de función de transferencia, su desempeño inferior no responde a un ajuste estadístico deficiente, sino a una limitación estructural intrínseca: la representación continua en términosde polos y ceros no incorpora términosde memoria sobre salidas pasadas en tiempo discreto, que son precisamente los que permiten al optimizador anticipar y corregir desviaciones acumulativasde larga duración propiasde las dinámicas lentasde la destilación. Este hallazgo sugiere que la funciónde transferencia identificable, pese a su popularidady simplicidadde implementación, no esla estructura más adecuada para columnasde destilación cuando el objetivo es el rechazode perturbacionesbajo MPC.

\vspace{0.25cm}

Finalmente, se concluye que el costo computacional no debe ser el criterio de selección del modelo en este tipo de aplicaciones. Dado que todos los modelos evaluados resuelven el problema de optimización en tiempos \'o rdenes de magnitud menores que el intervalo de muestreo, la factibilidad en tiempo real está garantizada para cualquier estructura. Esta constatación tiene una implicación práctica importante: el ingeniero de control no enfrenta un compromiso real entre precisión din\'a mica y viabilidad computacional en procesos con din\'a micas del orden de horas, por lo que puede y debe seleccionar el modelo atendiendo exclusivamente a su desempe\~n o en lazo cerrado.
